\section{Estimation and Inference}\label{sec:estimation}

\subsection{A menu-constrained propensity score}

Following \citet{angrist2011} and \citet{angrist2018}, the Committee's decision is modelled as an ordered probit
in a latent policy index $Z'\beta$ with cutoffs $c_1<c_2$, which gives unconstrained cell probabilities
$\tilde p_{-1}(z)=\Phi(c_1-z'\beta)$, $\tilde p_0(z)=\Phi(c_2-z'\beta)-\Phi(c_1-z'\beta)$ and
$\tilde p_1(z)=1-\Phi(c_2-z'\beta)$. To impose Assumption~\ref{as:menu} exactly, the probabilities are renormalised
over the menu:
\begin{equation}\label{eq:pscore}
  p_d(x,z)=\frac{\tilde p_d(z)\,\1\{d\in x\}}{\sum_{d'\in x}\tilde p_{d'}(z)} .
\end{equation}
This is the probability of choosing $d$ from an ordered probit whose outcome is known to lie in $x$. It is zero
exactly off the menu, it reduces to the ordered probit when $x=\cD$, and it is degenerate at one in the
single-option menus, which therefore do not enter the likelihood. The parameters $(\beta,c_1,c_2)$ are estimated
by maximum likelihood; standard errors come from the inverse Hessian, and average marginal effects use the
delta method.

The 2020 version entered the menus as dummy variables in an unconstrained ordered probit, which keeps the
probabilities off the menu small but positive; the renormalisation makes them exactly zero, as the bounds assume.

\paragraph{Covariates.} I consider two information sets. The \emph{macro} set follows \citet{angrist2018}: the
surprise component of federal funds futures prices, two lags of monthly PCE inflation, two lags of the change in
the unemployment rate, and two lags of year-on-year industrial-production growth. (The 2020 version used the
contemporaneous level of the IP index; growth rates lagged to the meeting date are used here instead.) The
\emph{Greenbook} set, which is the main specification, replaces the macro variables with the staff forecasts the
Committee had in hand at the meeting, as in \citet{romer2004}: current- and next-quarter forecasts of real GDP
growth and GDP-deflator inflation, and the current-quarter unemployment forecast, together with the futures
surprise. Each meeting is matched to the most recent Greenbook released before it.

Table~\ref{tab:pscore} reports average marginal effects on the probability of a rate increase. Columns (1)--(2)
reproduce the progression of the 2020 thesis on the meeting sample: adding the futures surprise raises the log
likelihood from $-110.0$ to $-91.4$. Column (3) imposes the menus through \eqref{eq:pscore}, with the same covariates
and no additional parameter, and the log likelihood rises by a similar amount again, to $-72.6$. The text of the
Bluebook carries as much information about the decision as market expectations do. Column (4), the main
specification, uses the Greenbook information set and fits best ($-52.7$). Higher current-quarter forecasts of
growth and inflation raise the probability of a hike. The next-quarter inflation forecast enters negatively once the
current-quarter forecast is held fixed; the two are correlated (0.63), so their separate coefficients are less
sharply determined than their sum. The futures surprise remains the strongest single predictor.

\begin{table}[!htbp]
  \centering
  \caption{Propensity-score models: average marginal effects on $\Prob(D_t=1)$}\label{tab:pscore}
  \small
  \tablefile{pscore.tex}
  \par\medskip
  \begin{minipage}{0.92\linewidth}\footnotesize
  \emph{Notes:} Ordered probit for the FOMC decision at 133 scheduled meetings, August 1989 -- January 2006.
  Columns (3)--(4) renormalise the probabilities over the Bluebook menu, equation~\eqref{eq:pscore}. Delta-method
  standard errors in parentheses. $^{*}$, $^{**}$, $^{***}$: significant at 10\%, 5\%, 1\%.
  GB = Greenbook forecast; $q_0$, $q_1$ = current and next quarter.
  \end{minipage}
\end{table}

\subsection{Estimating the bounds}

Given fitted scores $\hat p$, the conditional means in \eqref{eq:ipw} are estimated with normalised (H\'ajek)
inverse probability weights within each menu,
\[
\hat\mu_a(x)=\frac{\sum_{t:X_t=x}\1\{D_t=a\}\,Y_{t,h}/\hat p_a(W_t)}{\sum_{t:X_t=x}\1\{D_t=a\}/\hat p_a(W_t)},
\]
and $\pi_x$ by the sample share of meetings with menu $x$. Plugging these into Proposition~\ref{prop:bounds} gives
$\hat L_d$, $\hat U_d$ and $\hat\theta_d^{\cA}$. Normalised weights are bounded by the observed outcomes and are
less sensitive to small scores than the unnormalised estimator. In the main specification the smallest estimated
probability of an observed decision is 0.054, so no trimming is needed.

The outcome bounds $K_{1h},K_{2h}$ are set to the smallest and largest realised $h$-month change in each outcome
over August 1989 -- December 2010, a period that includes the 2001 and 2008--09 recessions. The assumption is that
the potential outcomes of any single meeting stay within the range the economy actually went through over these
two decades. The 2020 version calibrated $K$ to the interquartile range of the estimated effects across horizons;
calibrating it to the range of the outcome itself matches Assumption~\ref{as:bounded} and gives considerably wider
bounds.

\subsection{Inference}

The estimators are smooth functions of sample moments and of the probit estimates, but the observations are
serially dependent and the outcomes of nearby meetings overlap. I use a moving-block bootstrap over meetings
\citep{kunsch1989} with blocks of eight consecutive meetings (about one year) and $B=499$ draws. Each draw
re-estimates the propensity score, the menu shares and the conditional means, so first-stage uncertainty is
propagated. The 2020 version left inference for future work because the conditional effects are computed on
non-consecutive subsamples; resampling the full sequence of meetings and recomputing everything in each draw
addresses that concern.

For $\theta_d^{\cA}$ I report the basic bootstrap 90\% interval. For the partially identified $\theta_d$ I report the
\citet{imbens2004} interval, which covers the true parameter (rather than the whole identified set) with
probability 0.9 and remains valid when the identified set is narrow:
\[
\Bigl[\hat L_d-c_\alpha\,\hat\sigma_L,\ \hat U_d+c_\alpha\,\hat\sigma_U\Bigr],\qquad
\Phi\!\Bigl(c_\alpha+\tfrac{\hat U_d-\hat L_d}{\max(\hat\sigma_L,\hat\sigma_U)}\Bigr)-\Phi(-c_\alpha)=0.9,
\]
with $\hat\sigma_L,\hat\sigma_U$ the bootstrap standard errors. The support bounds $K$ are treated as known.
